\begin{lemma}
Let $T,J$ be finite and let $\nu$ be the nibble product measure with
$0\le p\le1$. Suppose $B_j$ are Borel events determined by finite sets
$S_j\subseteq T$: agreement on $S_j$ preserves membership in $B_j$.
Suppose each $S_j$ intersects at most $q$ other determining sets, where
$q\ge1$, and $\nu(B_j)\le P$ for a real $P\ge0$ satisfying
$eP(q+1)\le1$. Then $\nu(\bigcap_j B_j^c)>0$, and there exists an
outcome outside every $B_j$.
\end{lemma}
\begin{proof}
The factors of \noderef{NibbleProductMeasure} have mass one. A determining
event is a cylinder over its determining coordinates: fix the unused
coordinates at $(0,0)$ and take the inverse image of the event under this
measurable extension map. Its section is Borel and the determining property
identifies the original event with the cylinder over that section.
The product rectangle formula, extended to Borel sections by the pi-lambda
theorem, shows that an event is independent of the sigma algebra generated
by all events whose determining sets avoid its determining set. In particular
it is independent of the simultaneous avoidance of any of its nonneighbors.

Put $z=1/(q+1)$, so $0<z<1$.
Induct simultaneously on the size of a set $S$ of events to show that its
avoidance has positive probability and that, for $i\notin S$,
$\Pr[B_i\mid\bigcap_{j\in S}B_j^c]\le z$.
For the empty set positivity is one. For general $S$, its positivity
follows by ordering its members and using the already established bounds
for conditioning sets of smaller size, giving probability at least
$(1-z)^{|S|}>0$.
For the conditional bound, separate $S$ into neighbors $S_1$ of $i$ and
nonneighbors $S_2$. If $S_1$ is empty, independence gives a bound $P\le z$.
Otherwise the numerator conditioned on avoidance of $S_2$ is at most
$P$, by independence. In the denominator, order $S_1$ and expand its
avoidance conditional on $S_2$. Each factor is at least $1-z$ by induction:
its conditioning set has size at most $|S|-1$. Thus the desired conditional
probability is at most $P/(1-z)^{|S_1|}\le P/(1-z)^q$.
Since $q\log(1+1/q)\le1$, we have $(1-z)^q\ge e^{-1}$.
Consequently this ratio is at most $eP\le z$, completing the induction.
The avoidance probability of the entire family is at least
$(1-z)^{|J|}>0$. A positive-measure set is nonempty, giving the final
existence assertion. The argument also covers an empty event family.
\end{proof}
